\documentclass{subfiles}
\begin{document}
    Another attempt to what is todays cryptography, 
        is represented by substitution cyphers, for which we distinguish the mono-alphabetic ones,
        and the poly-alphabetic ones. 

    As the name suggests, mono-alphabetic substitution cyphers use a single alphabet to encypt the text.
        The core idea, implemented in several different ways (see \emph{AtBash} and \emph{Ceasar's cypher}),
        is that of replacing each character/symbol of the message with another of the same alphabet
        accordig to some rule. 
        For instance, Ceasar's cypher consists in replacing the occurrencies of the $n$-th character
        in the alphabet, with the $(n + k)$-th, for some $k$.
        AtBash itself can be seen as a Ceasar's cypher; in particular, one where $k = -2n$,
        i.e., $a$ is replaced with $z, b$ with $v$, etc.

    As it turns out, this kind of substitution cypher are weak.
        The reason of that comes comes from a simple statistical analysis.
        In particular, in each language, each character appears on average with a certain probability.
        Hence, if while analyzing a text we see a character that usually as low probability,
        appear very frequently, we can make a correspondence with the original symbol.
        For instance, in proper Italian, $j, k, w, x$ and $y$ have a zero probability to appear 
        in a text, thus, if we get a text where $w$ appears with a very high frequecy,
        most likely it replaces either an $n, p$ or a $t$ which are the most frequently used characters.

    Poly-alphabetic cyphers follow essentially the same principle, and are weak for the same reasons.
        Purely for historic reasons we cite \emph{Alberti's cyphering disk} and \emph{Vigen{\`e}re's cypher}.

\end{document}
